\documentclass[colorinlistoftodos]{article}
\usepackage[utf8]{inputenc}
\usepackage{natbib}
\usepackage[spanish]{babel}
\usepackage[numbib,nottoc]{tocbibind}
\usepackage{graphicx}
\usepackage{longtable}
\usepackage[hidelinks]{hyperref}
\usepackage{todonotes}


\title{
    {\large Universidad de Buenos Aires\vspace{20pt}}\\
    {\includegraphics[scale=0.4]{images/uba.eps}}\\
    {\large\textbf{Facultad de Ciencias Exactas y Naturales}}\vspace{16pt}\\
    {\small\textbf{Especialización en Exploración de Datos\\y Descubrimiento del Conocimiento}\vspace{16pt}}\\
    {\includegraphics[scale=0.5]{images/master.eps}}\\
    {Título\todo{Definir título}}
}
\author{Autora:\\Fernández Urquiza, Macarena}
\date{Fecha:\todo{Definir fecha}}

\begin{document}
\clearpage\maketitle
\thispagestyle{empty}

\newpage
\tableofcontents

\newpage

\begin{abstract}
\todo[color=blue!40]{Escribir abstract}Abstract.
\end{abstract}

\section{Introducción}

\subsection{Antecedentes}

% Comentario sobre el paper leído.
% Leer alguno de los otros papers descargados e
% introducir alguna reflexión sobre la importancia
% de la ingeniería de rasgos.

\subsection{Objetivos}

% objetivos grales: podrían ser los mismos que los
% del paper (?)
% objetivos particulares: podría ser contrastar si
% las observaciones realizada en el paper se
% verifican para el español

% GRAL
% Realizar una selección de rasgos que permita
% capturar las particularidades discursivas de los
% partidos a favor de la despenalización del aborto
% y de los partidos en contra

% PARTICULAR
% - Cotejar si las observaciones detalladas en el paper
% se evidencian también para el caso del español
% - Los rasgos seleccionados permiten entrenar un
% modelo que, dado un discurso, distinga a qué
% tipo de partido pertenece su emisor

\section{Datos}

% Fuente de los datos
% Procedimiento para su obtención (scrapper)
% Preprocesamiento realizado
% - .pdf --> .txt
% - pasaje a formato tabular
% Breve descripción de los datos:
% - cantidad de partidos totales
% - cantidad de oradores totales
% - distribución oradores~partidos
% - distribución género~partidos
% - palabras y tokens totales
% - distribución de palabras y tokens por partido
% - distribución de palabras y tokens por orador

\section{Metodología}

% Describir métodos aplicados.

%% Clasificación ==> w1: 0.004, w2: 0.002
%   - RF
%% Nonmodel-based ==> d(D-R) = w>0=D | w<0=R
%   - Dif. de frecuencias absolutas
%   - Dif. de proporciones
%   - Dif. de proporciones (s/stopwords)
%   - Odds
%   - Log-odds ratio
%   - Log-odds ratio (s/baja frecuencia)
%   - TF-IDF
%   - Word Scores
%% Model-based ???

%%% NaiveBayes (baseline)
%%% Reg. Logística

%       | RF | DFA | DFR |
% NB    |                |
% RL    |    |     |     |

\section{Resultados}

\section{Discusión y conclusiones}

\bibliographystyle{linquiry3}
\bibliography{bibliography}

\listoftodos%

\end{document}
